\documentclass{article}%
\usepackage[T1]{fontenc}%
\usepackage[utf8]{inputenc}%
\usepackage{lmodern}%
\usepackage{textcomp}%
\usepackage{lastpage}%
\usepackage{geometry}%
\geometry{margin=2.2cm}%
%
\usepackage{amsmath}%
\usepackage{amssymb}%
\usepackage{booktabs}%
\usepackage{longtable}%
\usepackage{array}%
\usepackage{hyperref}%
\hypersetup{hidelinks}%
\title{IME 2026 {-} Química (claude\_haiku\_5\_5) x Poliedro\\[4pt]\large Comparação com o gabarito Poliedro Resolve\thanks{Gerado com gpt-resolve: \url{https://github.com/lgabs/gpt-resolve}}}%
\date{10/11/2026}%
%
\begin{document}%
\normalsize%
\maketitle%
\section*{Resumo}\textbf{Nota média: 10,0/10} (16 de 16 itens com a resolução correta; cada item vale o mesmo dentro da questão).\par\medskip%
\begin{longtable}{@{}c c p{0.78\textwidth}@{}}\toprule Questão & Nota & Resumo \\ \midrule
1 & 10,0 & A resolução está certa e bate com o gabarito do Poliedro. Ela calcula Q = 9648 C = 0,1 F e n(M) = 0,05 mol (M$^2$? + 2e$^{-}$), o que dá M = 87,6 g/mol, logo o metal é o estrôncio (Sr). Conferi todas as contas, inclusive a tabela de candidatos (Rb, Y, Zr, Ca, Ba). Só há ressalvas pequenas nas verificações extras do item 7, que não afetam a resposta. \\
2 & 10,0 & O item a coincide: C é um álcool terciário. No item b, B (2{-}metilbut{-}1{-}eno) e C (2{-}metilbutan{-}2{-}ol) coincidem e estão corretos. Só A diverge: a resolução dá 1{-}cloro{-}2{-}metilbutano e o gabarito do Poliedro dá 2{-}cloro{-}2{-}metilbutano. Refazendo a análise, os dois haletos formam B por desidro{-}halogenação. O 1{-}cloro{-}2{-}metilbutano forma B como produto único, com qualquer base. O 2{-}cloro{-}2{-}metilbutano forma B como minoritário pela regra de Zaitsev, ou como majoritário com base volumosa. Como o enunciado não indica a base nem exige produto principal, as duas respostas são defensáveis. A da resolução é a mais rigorosa e deveria receber crédito integral. O gabarito apresenta uma única resposta sem ressalvas, e ela só vale sem restrição com base volumosa. \\
3 & 10,0 & A resolução está correta e bate com o gabarito do Poliedro nos dois itens. Em (a), as duas dão o acetato de linalila, CH2=CH--C(CH3)(OAc)--CH2CH2CH=C(CH3)2, com fórmula C12H20O2 (196,29 g/mol, conta conferida). Em (b), as duas dão a mesma unidade repetitiva {[}--C(CH3)2--CH(CH2CH2C(CH3)(OAc)CH=CH2)--{]}n, com a vinila pendente. Não há erro em nenhum dos dois. As dúvidas listadas são de baixa importância: a interpretação qualitativa de 'menor impedimento espacial' e o caráter hipotético da polimerização. \\
4 & 10,0 & A resolução está correta e coincide com o gabarito do Poliedro: são 120 mol de ácido acético, ou m $\approx$ 7,2 kg por piscina. A resolução usa M = 60,05 g/mol, tirada da tabela da prova, e chega a 7206 g. O Poliedro usa 60 g/mol e chega a 7200 g. As duas dão 7,2 kg, e a diferença é só de arredondamento. Os dois caminhos são equivalentes e corretos: a resolução usa balanço de carga mais Ka, e o gabarito usa Ostwald sem aproximação, com $\alpha$ = 1/6. Nenhum deles comete o erro de usar a aproximação {[}HAc{]} $\approx$ C, que daria 100 mol. Restam só hipóteses implícitas de baixa importância: ânion do sal neutro, ausência de alcalinidade na água e T = 25 $^\circ$C. \\
5 & 10,0 & A resolução e o gabarito Poliedro dão a mesma resposta, m1/m2 = 1 (m1 = m2), e os dois estão certos. Conferi o enunciado no PDF original e a transcrição está fiel. Refiz a conta: 324/81 = 4 e, pela proporção de 1 mol de CH4 para 2 mol de O2, m1/16 = 2$\cdot$m2/32, logo m1/m2 = 1. Com as massas da tabela da prova, a razão fica em $\approx$ 1,003, diferença que vem só do arredondamento das massas atômicas. As verificações extras da resolução estão corretas. A única hipótese implícita é que a segunda queima também é combustão de metano com água líquida; ela é razoável e a resolução a declara. \\
6 & 10,0 & A resolução está certa e coincide com o gabarito do Poliedro: ácido etanoico (acético), CH3COOH, H3C--C(=O)--OH, M = 60 g/mol, monoácido. As contas conferem: n(NaOH) = 3,0$\times$10$^{-}$$^3$ mol e M = 30 $\times$ 2 = 60 g/mol, o que dá k = 1 e a fórmula C2H4O2. As únicas diferenças estão no caminho: a resolução usa um argumento de desigualdade/associação mais complicado que o necessário e exclui explicitamente os isômeros que não são ácidos. Não há erro que afete a nota. \\
7 & 10,0 & A resolução está certa e coincide com o gabarito do Poliedro Resolve: 42 placas. Os valores intermediários também batem: T\_fus = 1296,48 K, $\Delta$H\_fus = 12,96 kJ/mol, $\Delta$H\_r = +24,0 kJ por mol de malaquita com H2O(g), E $\approx$ 83,82 MJ e N = 41,91. A resolução separa a energia em reação mais fusão, e o gabarito soma $\Delta$H\_fus dentro do $\Delta$H global; as duas formas são equivalentes. Há dois pontos de baixa importância para levar ao professor: a unidade de G no enunciado (kJ/mol) não combina com $\Delta$S\_fus = 10 J/(mol$\cdot$K), e o calor sensível foi desprezado por falta de dados de C\_p. Nenhum dos dois muda a resposta esperada. \\
8 & 10,0 & A resolução está certa e coincide com o gabarito Poliedro Resolve: $\approx$ 40,5 \%. O ânion tem 50 elétrons de valência (12 B $\times$ 3 + 12 H $\times$ 1 + 2 da carga). Um modelo só com ligações típicas de 2 centros e 2 elétrons precisaria de 84: 30 ligações B--B nas arestas do icosaedro mais 12 ligações B--H, ou seja, 42 ligações $\times$ 2. A deficiência é (84 - 50)/84 $\approx$ 40,5 \%. A verificação extra pela regra de Wade também está correta. Não encontrei erro nem na resolução nem no gabarito. \\
9 & 10,0 & Nos três itens, a resolução bate com o gabarito Poliedro: (a) 1,8$\times$10\^{}{-}5 mol/L, (b) pH $\approx$ 5,6 e (c) o pH diminui, porque a ionização é endotérmica ($\Delta$H$^\circ$ = +7,6 kJ/mol). Refiz todas as contas e elas estão corretas: $\Delta$G$^\circ$ = +36,3 kJ/mol, K $\approx$ 4,2$\times$10\^{}{-}7, $\Delta$H$^\circ$ = +7,6 kJ/mol. A única diferença é de modelo no item (b). O gabarito mantém o CO2(aq) fixo, como num sistema aberto, e chega a pH 5,56. A resolução desconta o CO2 ionizado com a equação do 2º grau e chega a 5,59. Os dois arredondam para 5,6. Mas a afirmação da resolução de que a equação do 2º grau é obrigatória é questionável, porque em equilíbrio com a atmosfera a abordagem do gabarito é a fisicamente correta. Esse ponto não altera a resposta final. \\
10 & 10,0 & A resolução está correta nos três itens e concorda com o gabarito Poliedro. (a) Ela traz a mesma sequência $\beta$$^{-}$, $\alpha$, $\alpha$, $\alpha$, identifica os nuclídeos (Ac, Th, Ra, Rn, Po) e acrescenta a equação global balanceada. (b) Classifica a série como 4n+3, do actínio, e justifica por que o resto da divisão por 4 não muda ao longo da série. (c) Dá $\lambda$ $\approx$ 0,17 s$^{-}$$^{1}$ e $\tau$ $\approx$ 5,8 s, que coincidem com 0,173 s$^{-}$$^{1}$ e 5,77 s do gabarito a menos de arredondamento. Não há erro de conta nem de conceito em nenhum dos dois. O único ponto a corrigir no gabarito é a nota de que o enunciado traria massa 277: no original consta 227. \\
\bottomrule\end{longtable}%
\section*{Principais pontos de dúvida}%
\begin{itemize}%
\item%
\textbf{Q2} \textit{(statement)} O enunciado não determina A de forma única. Tanto o 1{-}cloro{-}2{-}metilbutano (dá B como produto único) quanto o 2{-}cloro{-}2{-}metilbutano (dá B como minoritário com KOH/EtOH, ou como majoritário com t{-}BuOK) formam 2{-}metilbut{-}1{-}eno. O enunciado não indica a base nem diz que B é o produto principal. Vale confirmar com um professor, ou com o padrão oficial do IME, se as duas respostas para A são aceitas. Isso decide a nota do item b da resolução, que diverge do gabarito do Poliedro.%
\item%
\textbf{Q2} \textit{(key)} O gabarito do Poliedro escolhe A = 2{-}cloro{-}2{-}metilbutano sem ressalvas. Na desidro{-}halogenação com base comum (KOH alcoólico, regra de Zaitsev), o produto principal desse haleto é o 2{-}metilbut{-}2{-}eno, que na ozonólise dá acetona + acetaldeído e não formaldeído. B (2{-}metilbut{-}1{-}eno) só é majoritário com base volumosa (Hofmann). O gabarito não comenta isso nem cita o 1{-}cloro{-}2{-}metilbutano, que forma B exclusivamente.%
\item%
\textbf{Q9} \textit{(solution)} No item (b), a resolução diz que a aproximação z ? C não vale (15\% de ionização) e que a equação do 2º grau 'deve ser resolvida' com {[}CO2{]}eq = C - z. Isso trata C como concentração analítica de um sistema fechado. O item (a), porém, fala em equilíbrio com a atmosfera, e o item (c) também mantém o CO2 dissolvido constante. Num sistema aberto, a lei de Henry fixa o CO2(aq) livre em C, e o CO2 que se ioniza é reposto pelo ar. Assim, {[}H+{]} = $\surd$(K\_a1$\cdot$C) é exato (desprezando a água), e o pH dá 5,56. É o modelo do gabarito e a dedução clássica do pH 5,6 da chuva. O resultado final não muda (5,59 contra 5,56, ambos $\approx$ 5,6), mas a justificativa da resolução é conceitualmente questionável. Fica para o professor julgar se o enunciado admite a leitura de sistema fechado ('com base na concentração obtida\ldots{} solução resultante').%
\end{itemize}%
\section*{Questão 1}%
\href{https://poliedroresolve.sistemapoliedro.com.br/vestibulares/ime/2026/ime-2-fase-dia-3-25-10-2025-2-fase-dia-3/dissertativa-1-quimica-2-fase-dia-3-ime-2026}{Resolução do gabarito}\par\smallskip%
\begin{longtable}{@{}p{0.06\textwidth} p{0.33\textwidth} p{0.33\textwidth} p{0.1\textwidth} p{0.08\textwidth}@{}}\toprule Item & Resolução & Gabarito & Veredito & Certo \\ \midrule
Q1 -- identificação do metal M & Estrôncio (Sr, Z = 38): Q = 9648 C = 0,1 F; n\_M = 0,050 mol (z = 2); M\_m = 87,6 g/mol & Estrôncio (Sr): Q = 9648 C = 0,1 F; M$^2$? + 2e$^{-}$ $\to$ M; M\_M = 87,6 g/mol & igual & ambos \\
\bottomrule\end{longtable}%
\begin{itemize}%
\item%
\textbf{Q1 -- identificação do metal M)} Refiz as contas. t = 6,7 $\times$ 3600 = 24120 s; Q = 0,4 $\times$ 24120 = 9648 C = 0,1 F (com F = 96480 C/mol, como na folha de dados). Como z = 2, n\_M = 0,05 mol, e M = 4,38/0,05 = 87,6 g/mol. Na tabela da prova, o Sr tem 87,62 g/mol. A resolução e o gabarito chegam ao mesmo valor e à mesma conclusão. A tabela de conferência da resolução também está certa: Rb 8,55 g, Sr 4,381 g, Y 2,96 g, Zr 2,28 g, Ca 2,00 g, Ba 6,87 g. Os desvios também estão certos.%
\end{itemize}%
\paragraph{Pontos de dúvida}%
\begin{itemize}%
\item%
\textit{(solution)} No item 7 (rendimento), a resolução diz que, com $\eta$ = 0,90, M $\approx$ 97 g/mol e que 'não há metal da tabela com essa massa'. Isso é discutível: o Tc ($\approx$ 97--98) e o Mo (95,94) ficam bem perto. É uma verificação extra e não muda a resposta, que segue fixada pela hipótese de rendimento de 100\% e pela coincidência com o Sr.%
\item%
\textit{(solution)} A análise de z = 1 e z = 3 (43,8 e 131,4 g/mol) sobra, porque a fórmula MCl$_{2}$ já impõe z = 2. A afirmação de que não há metal nessa faixa é aproximada: o Sc (44,96) e o Cs (132,91) estão a cerca de 1--3\%. Não afeta a resposta.%
\item%
\textit{(statement)} O enunciado não diz que o rendimento faradaico é de 100\%. A resolução e o gabarito assumem isso implicitamente, que é a convenção usual. Com 0,4 A (um algarismo significativo), a identificação depende de tratar os dados como exatos, o que leva a Q = 0,1 F exato. Esse foi claramente o objetivo do examinador.%
\end{itemize}%
\section*{Questão 2}%
\href{https://poliedroresolve.sistemapoliedro.com.br/vestibulares/ime/2026/ime-2-fase-dia-3-25-10-2025-2-fase-dia-3/dissertativa-2-quimica-2-fase-dia-3-ime-2026}{Resolução do gabarito}\par\smallskip%
\begin{longtable}{@{}p{0.06\textwidth} p{0.33\textwidth} p{0.33\textwidth} p{0.1\textwidth} p{0.08\textwidth}@{}}\toprule Item & Resolução & Gabarito & Veredito & Certo \\ \midrule
a & Álcool (álcool terciário, 2{-}metilbutan{-}2{-}ol, sem H no carbono carbinólico, por isso não é oxidado pelo KMnO4/H+) & Álcool (álcool terciário, pois não reage com KMnO4 em meio ácido) & igual & ambos \\
b & A = 1{-}cloro{-}2{-}metilbutano, ClCH2--CH(CH3)--CH2--CH3; B = 2{-}metilbut{-}1{-}eno, CH2=C(CH3)--CH2CH3; C = 2{-}metilbutan{-}2{-}ol, (CH3)2C(OH)--CH2CH3 & A = 2{-}cloro{-}2{-}metilbutano, CH3CH2C(CH3)2Cl; B = 2{-}metilbut{-}1{-}eno, CH3CH2C(CH3)=CH2; C = 2{-}metilbutan{-}2{-}ol, CH3CH2C(CH3)2OH & diferente & ambos \\
\bottomrule\end{longtable}%
\begin{itemize}%
\item%
\textbf{a)} As duas respostas dão a mesma função orgânica (álcool) e a mesma justificativa: C é terciário e não tem H no carbono que leva o OH. A fórmula C5H10 + H2O $\to$ C5H12O, saturada, é coerente com isso.%
\item%
\textbf{b)} B e C coincidem. Só A diverge. Refiz a análise. A ozonólise que libera formaldeído exige CH2= terminal, o que deixa três candidatos: 1{-}penteno, 3{-}metil{-}1{-}buteno e 2{-}metil{-}1{-}buteno. Só o 2{-}metil{-}1{-}buteno dá exclusivamente álcool terciário: o cátion terciário formado não rearranja e fornece 2{-}metilbutan{-}2{-}ol. Logo B e C estão certos nas duas versões. Entre os 8 isômeros de C5H11Cl, dois formam 2{-}metilbut{-}1{-}eno. (1) O 1{-}cloro{-}2{-}metilbutano tem um único carbono $\beta$ com H (C2), então a E2 produz exclusivamente 2{-}metilbut{-}1{-}eno, com qualquer base. (2) O 2{-}cloro{-}2{-}metilbutano tem dois carbonos $\beta$ com H, C1 e C3. Com base pequena (KOH/EtOH, EtO-), o produto principal é o de Zaitsev, 2{-}metilbut{-}2{-}eno (\textasciitilde{}70\%), e o 2{-}metilbut{-}1{-}eno sai minoritário (\textasciitilde{}30\%). Com base volumosa (t{-}BuOK), o produto de Hofmann, 2{-}metilbut{-}1{-}eno, passa a ser o principal (\textasciitilde{}70\%). O enunciado diz apenas 'sofre desidro{-}halogenação formando B'. Ele não fala em produto principal nem indica a base. Por isso as duas escolhas de A são quimicamente defensáveis. A da resolução é mais robusta porque dá B como produto único sem nenhuma hipótese extra. A do gabarito só dá B como produto principal com base volumosa. Com a condição clássica (KOH alcoólico), B seria o produto minoritário. O gabarito não menciona essa ressalva nem a outra resposta possível. A resolução, por sua vez, exclui o 2{-}cloro{-}2{-}metilbutano usando uma hipótese própria (H1: 'formando B' = produto principal por Zaitsev) que não está no enunciado. Considero as duas respostas aceitáveis. Um corretor razoável deveria dar crédito integral à da resolução.%
\end{itemize}%
\paragraph{Pontos de dúvida}%
\begin{itemize}%
\item%
\textit{(statement)} O enunciado não determina A de forma única. Tanto o 1{-}cloro{-}2{-}metilbutano (dá B como produto único) quanto o 2{-}cloro{-}2{-}metilbutano (dá B como minoritário com KOH/EtOH, ou como majoritário com t{-}BuOK) formam 2{-}metilbut{-}1{-}eno. O enunciado não indica a base nem diz que B é o produto principal. Vale confirmar com um professor, ou com o padrão oficial do IME, se as duas respostas para A são aceitas. Isso decide a nota do item b da resolução, que diverge do gabarito do Poliedro.%
\item%
\textit{(key)} O gabarito do Poliedro escolhe A = 2{-}cloro{-}2{-}metilbutano sem ressalvas. Na desidro{-}halogenação com base comum (KOH alcoólico, regra de Zaitsev), o produto principal desse haleto é o 2{-}metilbut{-}2{-}eno, que na ozonólise dá acetona + acetaldeído e não formaldeído. B (2{-}metilbut{-}1{-}eno) só é majoritário com base volumosa (Hofmann). O gabarito não comenta isso nem cita o 1{-}cloro{-}2{-}metilbutano, que forma B exclusivamente.%
\item%
\textit{(solution)} A resolução exclui o 2{-}cloro{-}2{-}metilbutano pela hipótese H1 ('formando B' = produto principal por Zaitsev). Essa hipótese não está no enunciado. Na Observação, a resolução diz que a unicidade de A depende só de H5 (que exclui o 1{-}cloro{-}3{-}metilbutano). Na verdade ela também depende de H1: sem H1, ou com base volumosa, o 2{-}cloro{-}2{-}metilbutano também serve. A resposta final da resolução continua correta, mas a discussão de unicidade está incompleta.%
\item%
\textit{(solution)} Em H4, a resolução descreve o KMnO4 em meio ácido como 'diluído e frio'. Essa é a condição típica do teste de Baeyer, em meio neutro ou básico. O KMnO4/H+ costuma ser tratado como oxidante forte. A conclusão de que álcool terciário não é oxidado é a convenção usual em vestibular, mas a qualificação 'diluído e frio' é desnecessária e pouco precisa.%
\item%
\textit{(key)} Na imagem do gabarito, o produto da hidratação do 3{-}metilbut{-}1{-}eno (3{-}metilbutan{-}2{-}ol) aparece desenhado com uma dupla ligação junto ao carbono do OH. É erro de desenho e não afeta a resposta final. O gabarito também ignora o rearranjo 1,2{-}H, que na prática dá mistura de 3{-}metilbutan{-}2{-}ol e 2{-}metilbutan{-}2{-}ol. A exclusão se mantém, porque basta haver álcool secundário para o teste com KMnO4 dar positivo.%
\end{itemize}%
\section*{Questão 3}%
\href{https://poliedroresolve.sistemapoliedro.com.br/vestibulares/ime/2026/ime-2-fase-dia-3-25-10-2025-2-fase-dia-3/dissertativa-3-quimica-2-fase-dia-3-ime-2026}{Resolução do gabarito}\par\smallskip%
\begin{longtable}{@{}p{0.06\textwidth} p{0.33\textwidth} p{0.33\textwidth} p{0.1\textwidth} p{0.08\textwidth}@{}}\toprule Item & Resolução & Gabarito & Veredito & Certo \\ \midrule
a & CH2=CH--C(CH3)(OAc)--CH2--CH2--CH=C(CH3)2 (acetato de 3,7{-}dimetilocta{-}1,6{-}dien{-}3{-}ila), C12H20O2, 196,29 g/mol & CH2=CH--C(CH3)(OCOCH3)--CH2--CH2--CH=C(CH3)2 (acetato de 3,7{-}dimetilocta{-}1,6{-}dien{-}3{-}ila), C12H20O2, fórmula mínima C6H10O & igual & ambos \\
b & {[}--C(CH3)2--CH(CH2CH2C(CH3)(OAc)CH=CH2)--{]}n, por poliadição cabeça{-}cauda na dupla trissubstituída C6=C7, com a vinila mantida como grupo pendente & --{[}C(CH3)2--CH(CH2CH2C(CH3)(OCOCH3)CH=CH2){]}n--, por poliadição só na dupla C6=C7, com o vinil livre & igual & ambos \\
\bottomrule\end{longtable}%
\begin{itemize}%
\item%
\textbf{a)} Mesma estrutura e mesma fórmula molecular. Conferi as contas: C6H10O dá 98,14 g/mol e C12H20O2 dá 196,29 g/mol. A contagem de H fecha em 20 e o DBE é 3 (uma C=O e duas C=C). A cadeia mais longa tem 8 C (C1 a C8), as duplas ficam em 1 e 6 sem conjugação, e as duas metilas estão em C7. Metila e acetato vão em C3. A resolução justifica a escolha de C3 por tensão alílica A(1,3) com a metila cis em C7. O gabarito diz apenas que C3 fica 'mais afastado do grupo metil'. Os argumentos são diferentes, mas a conclusão é a mesma. Os dados extras da resolução (massa molar, C3 estereogênico, ausência de isomeria E/Z) estão corretos.%
\item%
\textbf{b)} A unidade repetitiva é a mesma nas duas respostas: C(CH3)2 e CH na cadeia principal, com o CH levando a cadeia lateral --CH2CH2C(CH3)(OAc)CH=CH2. Na resolução, o SMILES ({[}*{]}C(C)(C)C({[}*{]})CCC(C)(OC(C)=O)C=C) e o nome (poli{[}1,1{-}dimetil{-}2{-}(3{-}acetiloxi{-}3{-}metilpent{-}4{-}en{-}1{-}il)etileno{]}) estão coerentes com essa unidade. A regiosseletividade também está certa: a ligação nova se forma em C6 e o centro ativo terciário fica em C7. A unidade tem a fórmula do monômero, C12H20O2, com 196,29 g/mol.%
\end{itemize}%
\paragraph{Pontos de dúvida}%
\begin{itemize}%
\item%
\textit{(statement)} O critério 'carbono alílico, não terminal, de menor impedimento espacial' (v) é qualitativo, e os candidatos são C3 e C5. O gabarito escolhe C3 por ficar 'mais afastado do grupo metil'. A resolução chega a C3 pela tensão A(1,3) com a metila cis de C7. Os dois levam ao composto real (acetato de linalila), mas a banca pode aceitar justificativas diferentes.%
\item%
\textit{(solution)} No passo 2, a resolução diz que o terceiro grupo em C7 'só pode ser uma metila'. Isso só é obrigatório quando se combina com (iv). Pela estrutura sp2 de C7, ele também poderia ser H. É uma imprecisão de redação e não muda o resultado.%
\item%
\textit{(statement)} O enunciado chama o polímero de 'possivelmente formado'. Na prática, um alceno trissubstituído como C6=C7 homopolimeriza mal, e a via catiônica seria a mais plausível. Os dois lados tratam isso como exercício formal de poliadição e não especificam iniciador nem tacticidade. A resolução comenta esses pontos, mas eles não fazem parte da resposta esperada.%
\end{itemize}%
\section*{Questão 4}%
\href{https://poliedroresolve.sistemapoliedro.com.br/vestibulares/ime/2026/ime-2-fase-dia-3-25-10-2025-2-fase-dia-3/dissertativa-4-quimica-2-fase-dia-3-ime-2026}{Resolução do gabarito}\par\smallskip%
\begin{longtable}{@{}p{0.06\textwidth} p{0.33\textwidth} p{0.33\textwidth} p{0.1\textwidth} p{0.08\textwidth}@{}}\toprule Item & Resolução & Gabarito & Veredito & Certo \\ \midrule
Questão 4 -- massa mínima de ácido acético por piscina & C = 6,0$\times$10$^{-}$$^{4}$ mol/L; n = 120 mol; m $\approx$ 7,2 kg (120 $\times$ 60,05 g/mol $\approx$ 7206 g); pH final 4,00 & M = 6$\cdot$10$^{-}$$^{4}$ mol/L ($\alpha$ = 1/6); n = 120 mol; m = 120 $\times$ 60 = 7200 g = 7,2 kg & igual & ambos \\
\bottomrule\end{longtable}%
\begin{itemize}%
\item%
\textbf{Questão 4 -- massa mínima de ácido acético por piscina)} Refiz a conta. Com {[}H?{]} = 10$^{-}$$^{4}$ e Ka = 2,0$\times$10$^{-}$$^{5}$, {[}HAc{]} = h$^2$/Ka = 5$\times$10$^{-}$$^{4}$ e {[}Ac$^{-}${]} $\approx$ 10$^{-}$$^{4}$. Logo C = h(1 + h/Ka) = 6$\times$10$^{-}$$^{4}$ mol/L, o que corresponde a $\alpha$ = 1/6. Com V = 2$\times$10$^{5}$ L, isso dá 120 mol. A resolução e o gabarito chegam a esse valor por caminhos equivalentes: a resolução usa balanço de carga mais Ka, e o gabarito usa a Lei de Ostwald sem aproximação. A única diferença está na massa molar. A resolução usa 60,05 g/mol, tirada da tabela da prova (C 12,011; H 1,0079; O 15,999), e chega a 7206 g. O gabarito usa 60 g/mol e chega a 7200 g. As duas dão 7,2 kg, então a diferença é só de arredondamento. Conferi também os passos auxiliares da resolução: resolvendo o equilíbrio a partir de C = 6$\times$10$^{-}$$^{4}$, sai h = 1,0$\times$10$^{-}$$^{4}$; com C = 4$\times$10$^{-}$$^{4}$, sai h = 8$\times$10$^{-}$$^{5}$ (pH $\approx$ 4,10); o volume de ácido $\approx$ 6,9 L, uma fração de $\approx$ 3,4$\times$10$^{-}$$^{5}$ da piscina. Tudo está correto.%
\end{itemize}%
\paragraph{Pontos de dúvida}%
\begin{itemize}%
\item%
\textit{(statement)} O enunciado não diz se o ânion do sal contaminante é básico, nem se a água tem alcalinidade (HCO$_{3}$$^{-}$, CO$_{3}$$^2$$^{-}$). A resolução e o gabarito supõem, um explicitamente e o outro implicitamente, que todo H? fica livre em solução. Também supõem que a precipitação do metal não consome nem libera H?. A resolução declara essas hipóteses. O gabarito não as menciona.%
\item%
\textit{(statement)} A temperatura não é dada. Supor 25 $^\circ$C (Kw = 10$^{-}$$^{1}$$^{4}$, água neutra com pH 7) é razoável e não altera o resultado: a contribuição da água (10$^{-}$$^{1}$$^{0}$ mol/L de OH$^{-}$) é desprezível.%
\item%
\textit{(key)} O gabarito primeiro escreve Ka = M$\alpha$$^2$ com M$\alpha$ = 10$^{-}$$^{4}$, o que dá $\alpha$ = 0,2. Esse valor é só um teste de validade da aproximação, que é descartado em seguida pela forma completa Ka = M$\alpha$$^2$/(1-$\alpha$), com $\alpha$ = 1/6. A apresentação pode confundir, mas a conta final está correta. O gabarito também arredonda a massa molar para 60 g/mol, o que não muda os 7,2 kg.%
\item%
\textit{(solution)} A resolução interpreta a massa mínima como a que leva a pH = 4,00 exatamente, o limiar da condição pH $\leq$ 4, e justifica isso porque C(h) é monotônica. O gabarito adota a mesma interpretação. A resolução também usa uma densidade externa ($\approx$ 1,05 g/mL), mas só para uma verificação de volume, não no cálculo.%
\end{itemize}%
\section*{Questão 5}%
\href{https://poliedroresolve.sistemapoliedro.com.br/vestibulares/ime/2026/ime-2-fase-dia-3-25-10-2025-2-fase-dia-3/dissertativa-5-quimica-2-fase-dia-3-ime-2026}{Resolução do gabarito}\par\smallskip%
\begin{longtable}{@{}p{0.06\textwidth} p{0.33\textwidth} p{0.33\textwidth} p{0.1\textwidth} p{0.08\textwidth}@{}}\toprule Item & Resolução & Gabarito & Veredito & Certo \\ \midrule
Razão m1/m2 & m1/m2 = 1 (m1 = m2). Em geral, m1/m2 = M(CH4)/M(O), que dá $\approx$ 1,003 com as massas da tabela da prova. & m1/m2 = 1 (m1 = m2) & igual & ambos \\
\bottomrule\end{longtable}%
\begin{itemize}%
\item%
\textbf{Razão m1/m2)} Refiz a conta. Com E1/E2 = 324/81 = 4 e a proporção de 1 mol de CH4 para 2 mol de O2, n(CH4)\_1 = 4$\cdot$n(CH4)\_2 = 2$\cdot$n(O2)\_2. Daí m1/16 = 2$\cdot$m2/32, ou seja, m1/m2 = 1. Com C = 12,011, H = 1,0079 e O = 15,999, M(CH4)/M(O) = 16,043/15,999 $\approx$ 1,0027. A diferença vem só do arredondamento das massas atômicas. A resolução e o gabarito seguem o mesmo raciocínio: o gabarito usa x = 2y e a resolução cancela h. As contas de verificação da resolução também estão certas: 16$\cdot$324 = 32$\cdot$162 = 5184; $\Delta$Hc = -890,8 kJ/mol; m $\approx$ 5,8 g.%
\end{itemize}%
\paragraph{Pontos de dúvida}%
\begin{itemize}%
\item%
\textit{(statement)} O enunciado não diz explicitamente que a segunda queima também é de metano com água líquida. Diz só que foi feita "nas mesmas condições" e que consumiu oxigênio. A resolução e o gabarito admitem a mesma reação, CH4 + 2 O2 $\to$ CO2 + 2 H2O(l), e essa é a leitura natural. A resolução declara a hipótese; o gabarito a usa sem declarar.%
\item%
\textit{(solution)} A resolução dá a forma geral M(CH4)/M(O) $\approx$ 1,003 com as massas da tabela da prova e arredonda para 1. Isso é coerente com a precisão dos dados e não diverge do gabarito, que usa massas inteiras (16 e 32).%
\item%
\textit{(key)} A resolução do gabarito tem erros de digitação ou OCR que não mudam o resultado: "$\to$ \_ CO2", "calor liberado for de 81kJ" e a matriz de proporção mal formatada.%
\end{itemize}%
\section*{Questão 6}%
\href{https://poliedroresolve.sistemapoliedro.com.br/vestibulares/ime/2026/ime-2-fase-dia-3-25-10-2025-2-fase-dia-3/dissertativa-6-quimica-2-fase-dia-3-ime-2026}{Resolução do gabarito}\par\smallskip%
\begin{longtable}{@{}p{0.06\textwidth} p{0.33\textwidth} p{0.33\textwidth} p{0.1\textwidth} p{0.08\textwidth}@{}}\toprule Item & Resolução & Gabarito & Veredito & Certo \\ \midrule
Questão 6 -- nome e estrutura do ácido & Ácido etanoico (ácido acético), C2H4O2, CH3COOH, H3C--C(=O)--OH, M = 60 g/mol, monoácido (k = 1) & Ácido etanoico (ácido acético), C2H4O2 (M = 60 g/mol, monoácido): CH3--COOH, ou seja, H3C--C(=O)--OH & igual & ambos \\
\bottomrule\end{longtable}%
\begin{itemize}%
\item%
\textbf{Questão 6 -- nome e estrutura do ácido)} Conferi o enunciado no PDF (pág. 3): a transcrição está fiel. Contas refeitas: n(NaOH) = 0,1 $\times$ 0,030 = 3,0$\times$10$^{-}$$^3$ mol; pela densidade, M = 30 $\times$ 2 = 60 g/mol; pela titulação, k$\cdot$0,18/60 = 3,0$\times$10$^{-}$$^3$, o que dá k = 1 (monoácido). Com 12x + y + 32 = 60, sai x = 2 e y = 4: C2H4O2, ácido etanoico. A resolução segue outra ordem: primeiro tira M = 60k da titulação e depois limita M $\leq$ 60,5 com a densidade e um argumento de associação no vapor. Chega à mesma resposta e acrescenta a exclusão dos isômeros (HCOOCH3 e HOCH2CHO), que o gabarito não faz. Nome, fórmula molecular e estrutura plana são iguais nos dois.%
\end{itemize}%
\paragraph{Pontos de dúvida}%
\begin{itemize}%
\item%
\textit{(solution)} No passo 4, a resolução usa uma desigualdade (M\_ap $\geq$ M\_real, porque a associação só aumenta a massa molar aparente) e a massa do H igual a 1,0079 para concluir que 60k $\leq$ 60,5, logo k $\leq$ 1,008. O raciocínio está certo, mas é mais complicado do que precisa e pressupõe que o ácido não se decompõe no vapor. Além disso, a margem de 1\% é menor que a precisão dos dados (0,1 M e 0,18 g têm 1 a 2 algarismos significativos). O caminho direto do gabarito é mais robusto: M = 60 pela densidade e, então, k = 1 pela titulação.%
\item%
\textit{(solution)} Na conferência, a resolução escreve 60,05/2,016 $\approx$ 30. O valor exato é cerca de 29,8. É só arredondamento e não altera nada.%
\item%
\textit{(statement)} O enunciado idealiza o vapor como monomérico. Na realidade, o ácido acético dimeriza bastante no vapor, e a densidade medida perto da ebulição seria maior que 30 vezes a do H2. A expressão 'dissociação completa' deve ser lida como 'todos os H ionizáveis são neutralizados' (ou seja, contar os H+). Os dois textos interpretam assim, e a resolução observa que o ácido é fraco sem que isso mude a estequiometria.%
\end{itemize}%
\section*{Questão 7}%
\href{https://poliedroresolve.sistemapoliedro.com.br/vestibulares/ime/2026/ime-2-fase-dia-3-25-10-2025-2-fase-dia-3/dissertativa-7-quimica-2-fase-dia-3-ime-2026}{Resolução do gabarito}\par\smallskip%
\begin{longtable}{@{}p{0.06\textwidth} p{0.33\textwidth} p{0.33\textwidth} p{0.1\textwidth} p{0.08\textwidth}@{}}\toprule Item & Resolução & Gabarito & Veredito & Certo \\ \midrule
Questão 7 -- número mínimo de placas solares em 1 h & 42 placas (E $\approx$ 83,8 MJ; T\_fus = 1296,5 K; $\Delta$H\_fus = 12,96 kJ/mol; $\Delta$H\_r = +24,0 kJ por mol de malaquita; n\_Cu $\approx$ 3357,6 mol; N $\geq$ 41,9) & 42 placas ($\Delta$H global = +49,93 kJ por 2 mol de Cu; $\Delta$H $\approx$ 83,82 MJ; 83,82/2 = 41,91; T\_fus = 1296,48 K; $\Delta$H\_fus = 12,96 kJ/mol) & igual & ambos \\
\bottomrule\end{longtable}%
\begin{itemize}%
\item%
\textbf{Questão 7 -- número mínimo de placas solares em 1 h)} Refiz as contas. Pela diferença G\_L - G\_S = 12964,8 - 10,0T, T\_fus = 1296,48 K e $\Delta$H\_fus = 1296,48 $\times$ 10,0 = 12,96 kJ/mol. Com H2O(g), $\Delta$H\_r = (-242 - 788) - (-1054) = +24,0 kJ por mol de malaquita. Somando a fusão de 2 mol de Cu, a reação global dá $\Delta$H = +49,93 kJ por 2 mol de Cu. Com n\_Cu = 213360/63,546 = 3357,57 mol, a energia é E = 83,82 MJ e N = 41,91, logo N\_min = 42. Com M = 63,5 g/mol, E = 83,88 MJ e o resultado continua 42. A resolução separa a energia em reação (24,0/2 kJ por mol de Cu) mais fusão, e o gabarito põe H(Cu(l)) = +12,96 kJ dentro do $\Delta$H global. As duas formas são algebricamente iguais. A diferença entre 1296,5 e 1296,48 K é só arredondamento.%
\end{itemize}%
\paragraph{Pontos de dúvida}%
\begin{itemize}%
\item%
\textit{(statement)} O enunciado diz que G\_S e G\_L estão em kJ/mol. Mas a folha de dados dá $\Delta$S\_fus(Cu) = 10,0 J/(mol$\cdot$K), o que só é coerente se G estiver em J/mol: em kJ/mol, a inclinação de G\_L - G\_S daria 10 kJ/(mol$\cdot$K) e $\Delta$H\_fus $\approx$ 12965 kJ/mol. A resolução diz isso explicitamente e adota J/mol. O gabarito faz o mesmo sem comentar, pois usa $\Delta$S = 10 J/(mol$\cdot$K) e obtém $\Delta$H\_fus = 12,96 kJ. T\_fus não depende dessa escolha. Vale levar ao professor como possível erro de unidade no enunciado.%
\item%
\textit{(statement)} O termo '-2,7$\times$10$^2$' aparece nas duas expressões de G como uma constante solta. Provavelmente é erro de digitação de algo como -2,7$\times$10$^{-}$$^3$T$^2$. Isso não afeta o resultado, porque o termo é idêntico em G\_S e G\_L e se cancela na diferença.%
\item%
\textit{(solution)} Nem a resolução nem o gabarito contabilizam o calor sensível para aquecer reagentes e cobre até T\_fus. O enunciado não fornece C\_p, e a hipótese (iii) sugere usar $\Delta$H\_f constantes. A resolução discute isso e mostra que, somando só o aquecimento do Cu obtido de H\_S(T), chegaria a cerca de 81 placas. Esse caminho é incoerente com os dados e não é o pretendido pela banca, mas mostra que a resposta 42 depende do modelo adotado.%
\item%
\textit{(key)} O gabarito escreve 'H\_Cu(l) = +12,9648 kJ', ou seja, atribui ao Cu(l) uma 'entalpia de formação' igual a $\Delta$H\_fus. É um atalho de notação que só vale tomando H(Cu(s)) = 0 na mesma temperatura. Isso não muda o resultado. Ele também fala em 'células voltaicas' em vez de placas fotovoltaicas, o que é só redação.%
\item%
\textit{(solution)} A margem é pequena: N = 41,91, e a margem é de 0,2 MJ sobre 84 MJ. Ainda assim, para massas molares entre 63,5 e 63,55 g/mol, o arredondamento para cima dá sempre 42. Com T = 1358 K no lugar de T\_fus, daria 43, e a resolução avisa isso corretamente.%
\end{itemize}%
\section*{Questão 8}%
\href{https://poliedroresolve.sistemapoliedro.com.br/vestibulares/ime/2026/ime-2-fase-dia-3-25-10-2025-2-fase-dia-3/dissertativa-8-quimica-2-fase-dia-3-ime-2026}{Resolução do gabarito}\par\smallskip%
\begin{longtable}{@{}p{0.06\textwidth} p{0.33\textwidth} p{0.33\textwidth} p{0.1\textwidth} p{0.08\textwidth}@{}}\toprule Item & Resolução & Gabarito & Veredito & Certo \\ \midrule
Questão 8 -- deficiência percentual de elétrons de valência do (B12H12)2- & $\approx$ 40,5 \% (50 elétrons reais contra 84 esperados; diferença de 34; 34/84 $\times$ 100 = 850/21 $\approx$ 40,5 \%) & $\approx$ 40,5 \% (50 elétrons reais; 30 B--B + 12 B--H = 42 ligações $\times$ 2 = 84 esperados; (84 - 50)/84 $\times$ 100 $\approx$ 40,5 \%) & igual & ambos \\
\bottomrule\end{longtable}%
\begin{itemize}%
\item%
\textbf{Questão 8 -- deficiência percentual de elétrons de valência do (B12H12)2-)} A transcrição confere com o enunciado da página 4 do PDF. As duas contagens coincidem. Elétrons reais: 12$\cdot$3 + 12$\cdot$1 + 2 = 50. Arestas do icosaedro: o gabarito usa a relação de Euler, A = V + F - 2 = 30; a resolução usa 2A = 5V e confirma com Euler. Modelo 2c{-}2e: (30 + 12)$\cdot$2 = 84. Deficiência referida ao esperado, como pede o enunciado: 34/84 = 0,40476, ou seja, 40,5 \%. Conferi no python também os valores auxiliares da resolução: 16,7 \%, 42,9 \% e 68 \%. A verificação pela regra de Wade (26 = 2n + 2 elétrons de esqueleto para n = 12, cluster closo) está correta e é só complementar.%
\end{itemize}%
\paragraph{Pontos de dúvida}%
\begin{itemize}%
\item%
\textit{(statement)} O enunciado não define explicitamente a base do percentual, mas 'em relação ao esperado' indica o valor do modelo 2c{-}2e (84) como denominador. Tanto a resolução quanto o gabarito usam essa base. Com o valor real como base, o resultado seria 34/50 = 68 \%. A resolução comenta essa alternativa e a descarta.%
\item%
\textit{(solution)} O modelo hipotético 'só 2c{-}2e' dá 6 ligações a cada B (5 B--B e 1 B--H), ou seja, 12 elétrons em torno do boro, o que não é fisicamente realista. Mesmo assim, é o modelo que o enunciado pede, e a resolução e o gabarito o adotam igualmente. Isso não afeta a resposta.%
\end{itemize}%
\section*{Questão 9}%
\href{https://poliedroresolve.sistemapoliedro.com.br/vestibulares/ime/2026/ime-2-fase-dia-3-25-10-2025-2-fase-dia-3/dissertativa-9-quimica-2-fase-dia-3-ime-2026}{Resolução do gabarito}\par\smallskip%
\begin{longtable}{@{}p{0.06\textwidth} p{0.33\textwidth} p{0.33\textwidth} p{0.1\textwidth} p{0.08\textwidth}@{}}\toprule Item & Resolução & Gabarito & Veredito & Certo \\ \midrule
a & {[}CO2(aq){]} $\approx$ 1,8$\times$10\^{}{-}5 mol/L (P\_CO2 = 5,45$\times$10\^{}{-}4 atm; k\_H = 3,3$\times$10\^{}{-}2 mol/(atm$\cdot$L)) & {[}CO2{]} $\approx$ 1,8$\cdot$10\^{}{-}5 mol/L (S = k\_H$\cdot$P = 3,3$\cdot$10\^{}{-}2 $\cdot$ 545$\cdot$10\^{}{-}6) & igual & ambos \\
b & pH $\approx$ 5,6 (K\_a1 = e\^{}{-}14,68 $\approx$ 4,2$\times$10\^{}{-}7 a partir de $\Delta$G$^\circ$ = +36,3 kJ/mol; equação do 2º grau com {[}CO2{]}eq = C - z, {[}H+{]} $\approx$ 2,55$\times$10\^{}{-}6 mol/L, pH $\approx$ 5,59) & pH $\approx$ 5,6 ($\Delta$G$^\circ$ = +36,3 kJ/mol; K = e\^{}{-}14,68; {[}H+{]}$^2$ = 1,8$\cdot$10\^{}{-}5$\cdot$e\^{}{-}14,68; 2pH $\approx$ 11,11, pH $\approx$ 5,56) & igual & ambos \\
c & O pH diminui: $\Delta$H$^\circ$ = +7,6 kJ/mol (endotérmica), K\_a1 cresce com T e, com {[}CO2(aq){]} constante, {[}H+{]} aumenta ($\approx$ -0,02 de pH a cada 10 K) & O pH diminui: $\Delta$H$^\circ$ = +7,6 kJ/mol, a ionização é endotérmica, K e {[}H+{]} aumentam & igual & ambos \\
\bottomrule\end{longtable}%
\begin{itemize}%
\item%
\textbf{a)} Mesmo cálculo e mesmo valor: 3,3$\times$10\^{}{-}2 $\times$ 5,45$\times$10\^{}{-}4 = 1,7985$\times$10\^{}{-}5 mol/L. A observação da resolução sobre a correção de van't Hoff de 300 K para 298 K está certa (k\_H(298) $\approx$ 3,48$\times$10\^{}{-}2, cerca de 5\% maior; o pH mudaria só $\approx$0,01). É um complemento, não muda a resposta.%
\item%
\textbf{b)} $\Delta$G$^\circ$ = -586,8 - (-386,0 - 237,1) = +36,3 kJ/mol e ln K = -36300/(8,3$\cdot$298) = -14,676, logo K $\approx$ 4,23$\times$10\^{}{-}7. O gabarito mantém {[}CO2(aq){]} = C fixo e obtém pH = $\frac12$(pK + pC) $\approx$ 5,56. A resolução usa C - z no denominador e chega a pH $\approx$ 5,59. Os dois arredondam para 5,6. Há uma diferença de modelo, sem efeito no resultado: em equilíbrio com a atmosfera (sistema aberto), a lei de Henry fixa o CO2(aq) livre, e a abordagem do gabarito é a mais coerente fisicamente (é a dedução clássica do pH 5,6 da chuva). Ver doubts.%
\item%
\textbf{c)} $\Delta$H$^\circ$ = -692,0 - (-413,8 - 285,8) = +7,6 kJ/mol, confirmado. Pela equação de van't Hoff, K(308)/K(298) $\approx$ 1,105, então $\Delta$pH $\approx$ -$\frac12$$\cdot$log 1,105 $\approx$ -0,02. A estimativa da resolução está correta. Mesma conclusão e mesma justificativa.%
\end{itemize}%
\paragraph{Pontos de dúvida}%
\begin{itemize}%
\item%
\textit{(solution)} No item (b), a resolução diz que a aproximação z ? C não vale (15\% de ionização) e que a equação do 2º grau 'deve ser resolvida' com {[}CO2{]}eq = C - z. Isso trata C como concentração analítica de um sistema fechado. O item (a), porém, fala em equilíbrio com a atmosfera, e o item (c) também mantém o CO2 dissolvido constante. Num sistema aberto, a lei de Henry fixa o CO2(aq) livre em C, e o CO2 que se ioniza é reposto pelo ar. Assim, {[}H+{]} = $\surd$(K\_a1$\cdot$C) é exato (desprezando a água), e o pH dá 5,56. É o modelo do gabarito e a dedução clássica do pH 5,6 da chuva. O resultado final não muda (5,59 contra 5,56, ambos $\approx$ 5,6), mas a justificativa da resolução é conceitualmente questionável. Fica para o professor julgar se o enunciado admite a leitura de sistema fechado ('com base na concentração obtida\ldots{} solução resultante').%
\item%
\textit{(key)} Erros de digitação ou extração no gabarito: no item (c), aparece $\Delta$H\_f$^\circ${[}CO3\^{}{-}(aq){]} onde deveria ser HCO3\^{}{-}(aq), mas o valor usado (-692) é o do bicarbonato. No item (b), a equação aparece truncada ('H+ + HCO'). Nenhum dos dois afeta o resultado.%
\item%
\textit{(statement)} Hipóteses implícitas que tanto o gabarito quanto a resolução adotam: pressão total de 1 atm (ar seco, sem descontar a pressão de vapor d'água) e k\_H dada a 300 K usada a 25 $^\circ$C. A correção de van't Hoff ($\Delta$H\_sol = -19,8 kJ/mol) daria k\_H $\approx$ 3,48$\times$10\^{}{-}2 e mudaria o pH em apenas \textasciitilde{}0,01.%
\item%
\textit{(solution)} A resolução usa K\_w = 1,0$\times$10\^{}{-}14, valor que não está na folha de dados, para desprezar o {[}OH{-}{]}. É uma hipótese razoável e não afeta o resultado.%
\end{itemize}%
\section*{Questão 10}%
\href{https://poliedroresolve.sistemapoliedro.com.br/vestibulares/ime/2026/ime-2-fase-dia-3-25-10-2025-2-fase-dia-3/dissertativa-10-quimica-2-fase-dia-3-ime-2026}{Resolução do gabarito}\par\smallskip%
\begin{longtable}{@{}p{0.06\textwidth} p{0.33\textwidth} p{0.33\textwidth} p{0.1\textwidth} p{0.08\textwidth}@{}}\toprule Item & Resolução & Gabarito & Veredito & Certo \\ \midrule
a & Etapas: $^2$$^2$$^{7}$$_{8}$$_{9}$Ac $\to$ $^2$$^2$$^{7}$$_{9}$$_{0}$Th + $^{0}$?$_{1}$$\beta$$^{-}$ + $\nu$??; $^2$$^2$$^{7}$$_{9}$$_{0}$Th $\to$ $^2$$^2$$^3$$_{8}$$_{8}$Ra + $^{4}$$_{2}$$\alpha$; $^2$$^2$$^3$$_{8}$$_{8}$Ra $\to$ $^2$$^{1}$$^{9}$$_{8}$$_{6}$Rn + $^{4}$$_{2}$$\alpha$; $^2$$^{1}$$^{9}$$_{8}$$_{6}$Rn $\to$ $^2$$^{1}$$^{5}$$_{8}$$_{4}$Po + $^{4}$$_{2}$$\alpha$. Equação global: $^2$$^2$$^{7}$$_{8}$$_{9}$Ac $\to$ $^2$$^{1}$$^{5}$$_{8}$$_{4}$Po + 3 $^{4}$$_{2}$$\alpha$ + $^{0}$?$_{1}$$\beta$$^{-}$ + $\nu$??, com cargas e massas balanceadas. & $^2$$^2$$^{7}$$_{8}$$_{9}$X $\to$ $^2$$^2$$^{7}$$_{9}$$_{0}$J + $^{0}$?$_{1}$$\beta$; $^2$$^2$$^{7}$$_{9}$$_{0}$J $\to$ $^2$$^2$$^3$$_{8}$$_{8}$L + $^{4}$$_{2}$$\alpha$; $^2$$^2$$^3$$_{8}$$_{8}$L $\to$ $^2$$^{1}$$^{9}$$_{8}$$_{6}$Z + $^{4}$$_{2}$$\alpha$; $^2$$^{1}$$^{9}$$_{8}$$_{6}$Z $\to$ $^2$$^{1}$$^{5}$$_{8}$$_{4}$T + $^{4}$$_{2}$$\alpha$ & equivalente & ambos \\
b & Família 4n+3 (série do actínio / urânio{-}actínio): 227, 223, 219 e 215 deixam resto 3 na divisão por 4, e esse resto é invariante sob emissões $\alpha$ e $\beta$. & Série do Actínio (4n+3): os números de massa 227, 223, 219 e 215 deixam resto 3 na divisão por 4; X é o actínio. & igual & ambos \\
c & $\lambda$ = ln2/t$\frac12$ $\approx$ 0,694/4 $\approx$ 0,17 s$^{-}$$^{1}$; $\tau$ = 1/$\lambda$ = t$\frac12$/ln2 $\approx$ 5,8 s (unidades SI: s$^{-}$$^{1}$ e s). & $\lambda$ = (0,301/0,434)/4 $\approx$ 0,173 s$^{-}$$^{1}$; vida média $\tau$ = 1/$\lambda$ $\approx$ 5,77 s & equivalente & ambos \\
\bottomrule\end{longtable}%
\begin{itemize}%
\item%
\textbf{a)} As emissões de cada etapa coincidem: uma $\beta$$^{-}$ (A fixo, número atômico +1) e depois três $\alpha$ (A -4, número atômico -2). A resolução usa os símbolos reais (Ac, Th, Ra, Rn, Po), o que está correto pelos números atômicos. Ela também inclui o antineutrino e a soma global, com 227 = 215 + 12 e 89 = 84 + 6 - 1. Isso vai além do gabarito, mas não o contradiz. Conferi o enunciado original (p. 5): X e J têm massa 227, e não 277.%
\item%
\textbf{b)} Recalculei: 227, 223, 219 e 215 dão resto 3 na divisão por 4. A resolução ainda justifica por que o resto não muda ao longo da série, o que atende melhor ao verbo 'demonstre'.%
\item%
\textbf{c)} Refiz a conta: ln2 $\approx$ 0,301/0,434 = 0,6935, então $\lambda$ = 0,6935/4 = 0,1734 s$^{-}$$^{1}$ e $\tau$ = 1/$\lambda$ = 5,767 s. A resolução dá 0,17 s$^{-}$$^{1}$ e 5,8 s; o gabarito, 0,173 s$^{-}$$^{1}$ e 5,77 s. A diferença é só de arredondamento (dois algarismos contra três), e as unidades no SI estão corretas nos dois.%
\end{itemize}%
\paragraph{Pontos de dúvida}%
\begin{itemize}%
\item%
\textit{(key)} A nota do gabarito diz que o enunciado traz massa 277 para X e J. No PDF original (p. 5) aparece $^2$$^2$$^{7}$$_{8}$$_{9}$X e $^2$$^2$$^{7}$$_{9}$$_{0}$J, ou seja, 227. A observação do gabarito não procede, ao menos para esta versão da prova, e não afeta nenhum resultado.%
\item%
\textit{(solution)} A resolução apresenta $\lambda$ $\approx$ 0,17 s$^{-}$$^{1}$ e $\tau$ $\approx$ 5,8 s com dois algarismos significativos, embora t$\frac12$ seja dado como 'aproximadamente 4 s'. O gabarito usa três algarismos (0,173 s$^{-}$$^{1}$; 5,77 s). É só uma diferença de arredondamento e não deve mudar a nota.%
\item%
\textit{(statement)} O item (a) pede 'a equação estequiométrica da série' e não deixa claro se quer as etapas ou uma equação global. A resolução dá as duas formas, e o gabarito só as etapas. Incluir o antineutrino é fisicamente correto e opcional no nível da prova.%
\end{itemize}%
\end{document}